% =====================================================================
%  第 6 章 §6.6 —— Floer 同调与镜像对称（前半）
%  页段：原书 p.333–352（PDF 343–362）
%  译者：tr06
% =====================================================================

% [原文 p.333]
不过此式右端的含义很微妙，代入 $\hbar=0$ 它就变成常数。换言之，(6.67) 并不是在通常的绝对收敛意义下化为 (6.68) 的。更准确地说，(6.68) 的渐近展开才是 (6.67)。人们大概期待它与几何量子化、超局部解析有联系，但近来进展到何处，笔者还在学习中，不大清楚。除上面举出的之外，[36], [336], [388] 等也是形变量子化的基本文献。

Weinstein 在 [387] 中讨论了环面的情形，发现由 (6.68) 可把 $f\star_T g$ 表示为积分算子，且其积分核作为 $T$ 的函数满足非线性 Schrödinger 方程。环面情形的形变量子化，在从算子代数论走向非交换几何的进路（见 [62]）中很重要，已知它与非交换环面深刻相关（见 [313], [314], [386] 等）。[63] 等则讨论了非交换环面与所谓 M(atrix) 理论这一尝试\footnote{见 [30] 等。}（近来弦对偶研究的重要主题之一）的关系。

在形变量子化的研究中，辛拟群（symplectic groupoid）被热烈讨论，但笔者还没弄懂到能够解说的程度，故不叙述（见 [65], [336] 等）。

\section{Floer 同调与镜像对称}

\takeaway{同调镜像对称把 B 模型一侧的导出范畴（凝聚层、向量丛的 Ext 与 $A_\infty$ 结构）与 A 模型一侧的 Floer 同调（拉格朗日子流形与平坦 $U(1)$ 丛、Fukaya 范畴）对应起来；本节先造出 B 模型一侧的 Frobenius 代数与 $A_\infty$ 结构，再借助伪全纯圆盘把它翻译成 A 模型的语言。}

本节叙述同调镜像对称。先略微说明物理背景。§6.3 中已述：A 模型与 B 模型的真空的模空间，分别对应于辛结构与复结构的（扩大）模空间。1994–5 年前后开始的超弦理论第二次革命修正了这幅图景，

% [原文 p.334]
即：它们并不是 A 模型与 B 模型真空的全部，而只是“微扰真空”。“非微扰真空”的发现，是弦对偶研究的主要成果之一。

进而，经由其后的膜（brane）\footnote{膜的译名似乎还没有定论。笔者并非专家，考虑译名不是笔者之责。希望哪位专家来考虑。}理论，人们弄清了：“非微扰真空”对应于膜所处的某个状态。膜，朴素地说，弦是 1 维的线随时间运动所留下的痕迹，由 2 维流形表示；而膜一般是指维数高于 2 的子流形\footnote{关于膜的理论，参见 [303], [304] 等。}。理论的类型决定哪一维数的何种子流形作为膜出现。其机制在数学上尚未弄清\footnote{物理上的理解似乎在推进。}。这里不加说明地承认下述事项。

下面把这类暂时性的陈述称为“拟定义”“拟猜想”。

\begin{definition}{6.181}
  在由复流形 $(M,J_M)$ 紧化得到的 B 模型中，膜（brane）是指 $M$ 上解析凝聚层所成范畴的导出范畴的对象\footnote{这里不定义导出范畴（见 [173], [136] 等），本节也不深入到需要它的部分。复子流形、全纯向量丛、解析凝聚层等都给出解析凝聚层所成范畴的导出范畴的对象。}。

  在由辛流形 $(M,\omega_M)$ 紧化得到的 A 模型中，膜是指 $M$ 上的拉格朗日子流形 $L$ 与其上的平坦 $U(1)$ 丛 $\cL$ 所成的组 $(L,\cL)$。
\end{definition}

A 模型情形的拟定义 6.181 并不完整。尽可能使它接近准确的形式，是以下目标之一。由拟定义 6.181，在 A、B 两种模型中，膜的模空间大致上被确定了。即：B 模型中是 $(M,J_M)$ 上解析凝聚层所成范畴的导出范畴的对象的模空间；A 模型中是组 $(L,\cL)$ 的模空间。

\begin{conjecture}{6.182}
  当 $(M^\dagger,J_{M^\dagger})$ 是 $(M,\omega_M)$ 的镜像时，$(M^\dagger,J_{M^\dagger})$ 上的膜的模空间与 $(M,\omega_M)$ 上的膜的模空间一致。
\end{conjecture}

% [原文 p.335]
这就是 Kontsevich [215], [219] 所提出的同调镜像对称（homological mirror symmetry）猜想的一种形式\footnote{Kontsevich 提出同调镜像对称是在 1993 年前后，比超弦理论第二次革命开始得更早。Vafa [373] 出现之后，物理一侧似乎也开始研究同调镜像对称了。}。下面更详细地说明。

先考虑 B 模型。考虑导出范畴的对象会使讨论复杂，故考虑 $(M^\dagger,J_{M^\dagger})$ 上的全纯向量丛 $E$。如 §6.4 所述，全纯向量丛模空间在 $E$ 附近的邻域由仓西映射 $H^1_{\bar\partial}(M;\mathrm{Hom}(E,E))\to H^2_{\bar\partial}(M;\mathrm{Hom}(E,E))$ 的零点给出。若考虑扩大模空间，并把模空间看作超流形，则扩大模空间的“切空间”是 $H^*_{\bar\partial}(M;\mathrm{Hom}(E,E))=\mathrm{Ext}^*(E,E)$。（用后一记号，对解析凝聚层及其所成范畴的导出范畴的对象也同样成立。）没有膜的情形，即考虑复流形模空间时，复结构扩大模空间的切空间是 $H^*_{\bar\partial}(M;\Omega_{\C})$。§6.5 所作的构成，是在复结构扩大模空间上定义 Frobenius 流形结构。我们想对膜的模空间作类似的尝试，即在模空间的切空间上定义 Frobenius 代数结构。

设 $M$ 为 Calabi–Yau 复流形，$E$ 为其上秩 $m$ 的全纯向量丛。设 $\Omega$ 为 $M$ 上的全纯 $n$ 次形式。令 $\mathfrak{H}^*_B(M;E)=\mathrm{Ext}^*(E,E)$。把积 $\mathfrak{m}^B_2\colon\mathfrak{H}^{k_1}_B(M;E)\otimes\mathfrak{H}^{k_2}_B(M;E)\to\mathfrak{H}^{k_1+k_2}_B(M;E)$ 取作米田积。向量丛映射 $\mathrm{Tr}\colon\mathrm{Hom}(E,E)\to\C$ 由在各点取迹来定义。$\mathrm{int}_M\colon\mathfrak{H}^n_B(M;E)\to\C$ 定义为
\[
  \mathrm{int}_M(\varphi\wedge u)=\int_M\mathrm{Tr}(\varphi)u\wedge\Omega
\]
（$\varphi\in\Gamma(M;\mathrm{Hom}(E,E))$, $u\in\Gamma(M;A^{0,*}M)$）。内积 $\langle\cdot,\cdot\rangle_B$ 由下

% [原文 p.336]
式定义：
\[
  \langle u_1,u_2\rangle_B=\mathrm{int}_M\bigl(\mathfrak{m}^B_2(u_1,u_2)\bigr).
\]

\begin{lemma}{6.183}
  $(\mathfrak{H}^*_B(M;E),\mathfrak{m}^B_2,\langle\cdot,\cdot\rangle_B)$ 是 Frobenius 代数。
\end{lemma}

证明容易。

\begin{remark}{6.184}
  即使假设 $M$ 是 Calabi–Yau 流形，其上的向量丛模空间所对应的仓西映射也未必是 0。即定理 6.142 的类似并不成立。因此，膜的模空间一般不是 Frobenius 流形。
\end{remark}

进而还可定义（高次）Massey–米田积。积决定仓西映射的 2 次项，Massey–米田积决定 3 次项。（高次）Massey–米田积合起来定出 $A_{\infty}$ 结构。本节的主要目的，就是在 A 模型的情形下构造这个 $A_{\infty}$ 结构的镜像。

\begin{remark}{6.185}
  在转入 A 模型的讨论之前，关于 B 模型中的膜的模空间再稍作叙述。

  全纯向量丛与解析凝聚层的模空间，在代数几何中已被许多人研究。它也被应用于用 Donaldson 不变量研究 4 维微分流形。近来出现了与本节内容相关的、下面略述的发展。

  Donaldson–Thomas [80], Thomas [364] 猜想：对复 3 维 Calabi–Yau 流形及其上的复向量丛 $E$，$E$ 上全纯向量丛结构的个数会成为 Calabi–Yau 流形的不变量\footnote{Minahan–Nemeschansky–Vafa–Warner 的论文 [260] 似乎与这个猜想有关。}。（这个模空间的虚拟维数是 0，故数其点数有意义。）Donaldson–Thomas 提出的不变量称为复 Casson 不变量（holomorphic Casson invariant）\footnote{Casson 不变量曾是这样一个不变量：对 3 维流形上的向量丛，对应其“平坦联络的个数”。把平坦联络换成自对偶联络，把实 3 维换成复 3 维，普通的 Casson 不变量就变成复 Casson 不变量。关于 Casson 不变量，见 [5]。}。
\end{remark}

% [原文 p.337]
进而 Donaldson–Thomas 通过考察退化的极限，找到了可计算复 Casson 不变量的例子。这个计算法与用 Heegaard 分解（Heegaard splitting）计算 Casson 不变量的方法类似。

要证明复 Casson 不变量是真正的不变量，模空间的紧化研究不可或缺。为此，人们认为定义 Gromov–Witten 不变量时所用的方法很有力，Thomas 等人正在推进其完成。

另一方面，Tian [366] 研究了模空间的紧化，其中包含全纯向量丛的模空间与 4 维流形上自对偶联络的模空间\footnote{Donaldson 理论中所用者。}，属于更一般的“自对偶联络”情形\footnote{具有例外和乐的黎曼流形，即在 7 维以 $G_2$ 为和乐群的情形、在 8 维以 $\mathrm{Spin}(7)$ 为和乐群的情形，是另外的重要例子。它们称为 Joyce 流形（Joyce manifold）。Joyce 流形由 [52], [193], [194] 构成，其上的规范场论见 [57] 等。}。联络空间的紧化，必须研究曲率发散（爆炸）的点集的情形。Nakajima [275], [276] 用 Uhlenbeck 的方法，对一般维数的 Yang–Mills 联络序列证明了曲率爆炸点集的余维数在 Hausdorff 维数（Hausdorff dimension）意义下为 4。Tian 证明了曲率爆炸点集是可求长（rectifiable）\footnote{意谓在几乎所有点处切空间存在。见几何测度论（geometric measure theory）的教科书 [98], [172], [339]。}的，并进一步证明了它在注意 2.84 的意义下是校准的“子流形”\footnote{一般带有奇点。}。

以上研究可以说是高维规范场论的一部分。规范场论的数学研究长期几乎只限于 4 维情形，但高维情形的研究也在兴盛起来。

复流形上全纯向量丛的模空间，是 B 模型中膜的模空间的重要部分；另一重要部分是复子流形的模空间，即 Hilbert 概形（Hilbert scheme）。以 0 维复子流形的模空间、即点 Hilbert 概形为中心，中岛等人的研究正在推进（参见 [280]）。

% [原文 p.338]
结合后面要讲的猜想 6.193 来看，成为（特异）环面纤维空间的复流形上，全纯向量丛模空间的研究很重要。[111], [112], [113], [71] 等做了这项研究\footnote{不过所研究的似限于纤维为复环面的情形。与同调镜像对称相关的是纤维为实拉格朗日子流形的情形。[111], [112], [113], [71] 等所用的代数几何手法对实环面纤维是否有效，笔者不大清楚。}。

Simpson [340] 用 $n$ 范畴（$n$ category）等展开了以非交换群为系数的高次上同调的 Hodge 理论，这也许与本节的构成有某种关联。

在 A 模型中，拉格朗日子流形 $L$ 与其上的复 $U(1)$ 丛 $\cL$ 所成组 $(L,\cL)$ 的模空间 $\mathcal{LAG}_{\mathrm{class}}(M)$，对应于全纯向量丛、解析凝聚层等的模空间——这便是拟猜想 6.182。（$\mathcal{LAG}_{\mathrm{class}}(M)$ 的完全严格的定义尚未给出。后面再说明。）

如注意 6.13 所述，从复流形出发的 B 模型中没有量子修正。这件事的一个推论是：全纯向量丛、解析凝聚层等的模空间本身，在 B 模型中就成为膜的（量子化）模空间。

另一方面，A 模型中存在量子修正。即 A 模型的膜的（量子化）模空间不是 $\mathcal{LAG}_{\mathrm{class}}(M)$ 本身，而是它的量子变形 $\mathcal{LAG}_{\mathrm{qm}}(M)$。因此，当辛流形 $M$ 的镜像是复流形 $M^\dagger$ 时，$\mathcal{LAG}_{\mathrm{qm}}(M)$ 应与 $M^\dagger$ 上全纯向量丛、解析凝聚层等的模空间一致——这是同调镜像对称猜想的一种形式。

研究这一点，就能把模空间的量子变形这一神秘现象作为严格的数学来研究，这是其魅力之一。

先记述 $\mathcal{LAG}_{\mathrm{class}}(M)$。辛流形 $M$ 中，拉格朗日子流形全体构成无限维的族。因此，要得到有限维的模空间，必须加上某种条件，或者考虑等价关系。